%Version 3.1 December 2024
% See section 11 of the User Manual for version history
%
%%%%%%%%%%%%%%%%%%%%%%%%%%%%%%%%%%%%%%%%%%%%%%%%%%%%%%%%%%%%%%%%%%%%%%
%%                                                                 %%
%% Please do not use \input{...} to include other tex files.       %%
%% Submit your LaTeX manuscript as one .tex document.              %%
%%                                                                 %%
%% All additional figures and files should be attached             %%
%% separately and not embedded in the \TeX\ document itself.       %%
%%                                                                 %%
%%%%%%%%%%%%%%%%%%%%%%%%%%%%%%%%%%%%%%%%%%%%%%%%%%%%%%%%%%%%%%%%%%%%%

%%\documentclass[referee,sn-basic]{sn-jnl}% referee option is meant for double line spacing

%%=======================================================%%
%% to print line numbers in the margin use lineno option %%
%%=======================================================%%

%%\documentclass[lineno,pdflatex,sn-basic]{sn-jnl}% Basic Springer Nature Reference Style/Chemistry Reference Style

%%=========================================================================================%%
%% the documentclass is set to pdflatex as default. You can delete it if not appropriate.  %%
%%=========================================================================================%%

%%\documentclass[sn-basic]{sn-jnl}% Basic Springer Nature Reference Style/Chemistry Reference Style

%%Note: the following reference styles support Namedate and Numbered referencing. By default the style follows the most common style. To switch between the options you can add or remove “Numbered” in the optional parenthesis. 
%%The option is available for: sn-basic.bst, sn-chicago.bst%  
 
%%\documentclass[pdflatex,sn-nature]{sn-jnl}% Style for submissions to Nature Portfolio journals
%%\documentclass[pdflatex,sn-basic]{sn-jnl}% Basic Springer Nature Reference Style/Chemistry Reference Style
\documentclass[pdflatex,sn-mathphys-num]{sn-jnl}% Math and Physical Sciences Numbered Reference Style
%%\documentclass[pdflatex,sn-mathphys-ay]{sn-jnl}% Math and Physical Sciences Author Year Reference Style
%%\documentclass[pdflatex,sn-aps]{sn-jnl}% American Physical Society (APS) Reference Style
%%\documentclass[pdflatex,sn-vancouver-num]{sn-jnl}% Vancouver Numbered Reference Style
%%\documentclass[pdflatex,sn-vancouver-ay]{sn-jnl}% Vancouver Author Year Reference Style
%%\documentclass[pdflatex,sn-apa]{sn-jnl}% APA Reference Style
%%\documentclass[pdflatex,sn-chicago]{sn-jnl}% Chicago-based Humanities Reference Style

%%%% Standard Packages
%%<additional latex packages if required can be included here>

\usepackage{graphicx}%
\usepackage{multirow}%
\usepackage{amsmath,amssymb,amsfonts}%
\usepackage{amsthm}%
\usepackage{mathrsfs}%
\usepackage[title]{appendix}%
\usepackage{xcolor}%
\usepackage{textcomp}%
\usepackage{manyfoot}%
\usepackage{booktabs}%
\usepackage{algorithm}%
\usepackage{algorithmicx}%
\usepackage{algpseudocode}%
\usepackage{listings}%
%%%%

%%%%%=============================================================================%%%%
%%%%  Remarks: This template is provided to aid authors with the preparation
%%%%  of original research articles intended for submission to journals published 
%%%%  by Springer Nature. The guidance has been prepared in partnership with 
%%%%  production teams to conform to Springer Nature technical requirements. 
%%%%  Editorial and presentation requirements differ among journal portfolios and 
%%%%  research disciplines. You may find sections in this template are irrelevant 
%%%%  to your work and are empowered to omit any such section if allowed by the 
%%%%  journal you intend to submit to. The submission guidelines and policies 
%%%%  of the journal take precedence. A detailed User Manual is available in the 
%%%%  template package for technical guidance.
%%%%%=============================================================================%%%%

%% as per the requirement new theorem styles can be included as shown below
\theoremstyle{thmstyleone}%
\newtheorem{theorem}{Theorem}%  meant for continuous numbers
%%\newtheorem{theorem}{Theorem}[section]% meant for sectionwise numbers
%% optional argument [theorem] produces theorem numbering sequence instead of independent numbers for Proposition
\newtheorem{proposition}[theorem]{Proposition}% 
%%\newtheorem{proposition}{Proposition}% to get separate numbers for theorem and proposition etc.

\theoremstyle{thmstyletwo}%
\newtheorem{example}{Example}%
\newtheorem{remark}{Remark}%

\theoremstyle{thmstylethree}%
\newtheorem{definition}{Definition}%

\raggedbottom
%%\unnumbered% uncomment this for unnumbered level heads

\begin{document}

\title[Article Title]{Inter-Subject Correlation of Heart Rate and Attention Modulation}

%%=============================================================%%
%% GivenName	-> \fnm{Joergen W.}
%% Particle	-> \spfx{van der} -> surname prefix
%% FamilyName	-> \sur{Ploeg}
%% Suffix	-> \sfx{IV}
%% \author*[1,2]{\fnm{Joergen W.} \spfx{van der} \sur{Ploeg} 
%%  \sfx{IV}}\email{iauthor@gmail.com}
%%=============================================================%%

\author{\fnm{Roberto} \sur{Antoniello 47891A}} 

%%==================================%%
%% Sample for unstructured abstract %%
%%==================================%%

\abstract{ This project focuses on analyzing heart-rate synchrony across individuals to verify if it's modulated by attentional state while watching educational videos.

The goal is therefore to replicate a couple of results from Pérez et al. (2021)~\cite{perez2021} using data from a newer dataset: the Brain, Body and Behavior Dataset (BBBD)~\cite{bbbd}.

To achieve this, heart-rate fluctuations were extracted from raw ECG recordings and compared across subjects using ISC-HR (inter-subject correlation of heart-rate).

A hypothesis test with a FDR correction is used to establish statistical significance of the obtained results. ISC-HR is then compared between attentive and distracted conditions through a paired statistical test across all five stimuli.

Hypothesis test and descriptive statistics are employed to analyze synchrony patterns and to detect the effect of attention across stimuli.

The results provide a statistically grounded replication of ISC-HR modulation by attention, discussed in relation to the original paper's finding on which this project is based.

 }

%%\pacs[JEL Classification]{D8, H51}

%%\pacs[MSC Classification]{35A01, 65L10, 65L12, 65L20, 65L70}

\maketitle

\section{Introduction}\label{sec1}

Heart rate is often treated as a purely autonomic signal, reflecting internal physiological states such as arousal or stress.\\
Recent studies suggest that heart-rate fluctuations can be driven by external stimuli.
When multiple individuals are exposed to the same narrative content, their heart rates tend to fluctuate in a similar way.\\
This phenomenon is known as inter-subject correlation of heart-rate (ISC-HR).\\

Pérez et al. (2021)~\cite{perez2021} proposed that these heart-rate fluctuations are modulated by conscious, attentive processing of the stimulus. \\

This project aims to replicate two of the main results reported by the analyzed paper.
\begin{itemize}
\item ISC-HR is statistically significant across subjects exposed to the same stimulus (R1).
\item Synchrony is modulated by attentional state, being higher when subjects are attentive than when they are distracted (R2).
\end{itemize}


This work uses data from the Brain, Body and Behavior Dataset (BBBD)~\cite{bbbd}, which is a more recent dataset that includes similar experiments to the original study.

\section{Related Work}\label{Sec2}

The original paper~\cite{perez2021} computes ISC-HR by correlating each subject's instantaneous heart-rate fluctuations against the rest of the group, averaging the resulting pairwise correlations in Fisher-Z space, and testing significance via a circular-shift permutation procedure with False Discovery Rate (FDR) correction.
The study reports three main findings across four experiments.\\


The first one (R1) shows that heart-rate fluctuations synchronize across subjects listening to the same auditory narrative, with 17 of 27 subjects having ISC-HR statistically significant.\\

The second result (R2) shows this synchrony is systematically higher when subjects attend to a stimulus than when they are distracted.\\

The third result (R3) shows that ISC-HR predicts subsequent memory recall performance for the narrative content, suggesting that the measure reflects conscious engagement rather than an autonomic response to sensory input.\\

The fourth experiment applies ISC-HR to patients with disorders of consciousness, showing reduced synchrony relative to healthy controls and framing ISC-HR as a potential prognostic marker in a clinical setting.

As already anticipated, this project replicates only R1 and R2 using an experiment equivalent to the original study's experiment 2.\\

\section{Dataset}\label{sec3}

This project uses Experiment 2 of the BBBD dataset~\cite{bbbd}, which contains recordings of ECG, EEG, eye-tracking, and respiration during educational video viewing.\\

In experiment 2, 31 subjects watch five educational videos, each one viewed twice by the same subject: attentively the first time, distracted by performing a secondary task (silently counting backward) the second time.

BBBD's data descriptor~\cite{bbbd} explicitly states that the ECG data provided for this experiment are not newly collected, but are the same recordings originally published in Pérez et al., 2021~\cite{perez2021}.\\

ECG preprocessing documentation~\cite{bbbd} states that R-wave peaks were detected automatically and then manually corrected.
This project instead detects R-peaks directly from the raw ECG signal using an automated procedure, without using BBBD's provided peak timestamps and without manual review.\\
This choice was made in order to reproduce and understand the original study's methodology.\\

Not every subject has a usable recording for every (subject, session, stimulus) combination. Missing files are skipped automatically, so the number of subjects analyzed varies slightly across stimuli and conditions.


\section{Methodology}\label{sec4}

introduction to methods used
Qui hai già quasi tutto pronto, solo da sintetizzare (i documenti supplementari sono la versione "lunga", qui ne fai il riassunto):


\subsection{Signal preprocessing}\label{subsec1}

da HR-Pipeline-Note.pdf + Butterworth-Filter-Note.pdf: filtro, picchi, HR istantanea, interpolazione

\subsection{ISC-HR computation}\label{subsec2}

da Hypothesis-Test-R1-Note.pdf Sezione 1: Fisher-Z, Step 3-5 del paper


\subsection{Hypothesis test}\label{subsec3}

permutazione circolare, Benjamini-Hochberg


\subsection{Attention modulation}\label{subsec4}

da R2-Attention-Modulation-Note.pdf: correlazione a riferimento fisso, esclusione dell'auto-correlazione, Shapiro-Wilk + t-test/Wilcoxon.

\subsection{Implementation choices}\label{subsec5}

(opzionale, se vuoi una tabella riassuntiva) — QC, i parametri non specificati dal paper, tutti già tracciati nei documenti

\section{Results}\label{sec5}

Qui attingi principalmente da Final-Results-R1-R2-Discussion.pdf, che è già scritto quasi in forma di risultati/discussione:

\subsection{Input validation}\label{subsec6}

\subsection{R1 on all combinations of (session, stimulus)}\label{subsec7}

\subsection{Significant stimulus}\label{subsec8}

\subsection{Stimulus 02 case}\label{subsec9}


\section{Discussion}\label{sec6}

Confronto con le cifre del paper (con le cautele che abbiamo già scritto), limiti dichiarati (niente ANOVA completa, sensibilità all'interpolazione non testata, scelte QC/esclusione non specificate dal paper), interpretazione del pattern generale.

\section{Conclusion}\label{sec7}

Cosa hai dimostrato, cosa resta aperto (R3, metodo esatto per test multipli, ecc.).

8. References

Paper Pérez et al., documentazione BBBD/Zenodo, eventuali librerie citate.




\section*{Declarations}
\textbf{Ai Usage Disclaimer:} parts of this project have been realized with the assistance of ChatGPT(model GPT-5.3) and Claude(model Sonnet 4.6) for the following types of tasks:
\begin{itemize}
\item Reviewing and refining text during the writing of this report.
\item Collecting ideas about results analysis.
\item Generation of Python code to support the development.
\end{itemize}

Every content generated by AI has been reviewed and tested by the author.


\begin{thebibliography}{99}

\bibitem{perez2021}
P.~P{\'e}rez, J.~Madsen, L.~Banellis, B.~T{\"u}rker, F.~Raimondo, V.~Perlbarg, M.~Valente, M.-C.~Ni{\'e}rat, L.~Puybasset, L.~Naccache, T.~Similowski, D.~Cruse, L.~C.~Parra, and J.~D.~Sitt.
\newblock Conscious processing of narrative stimuli synchronizes heart rate between individuals.
\newblock \emph{Cell Reports}, 36(11):109692, 2021.

\bibitem{bbbd}
J.~Madsen, N.~Kuppa, and L.~C.~Parra.
\newblock The Brain, Body, and Behavior Dataset (BBBD): Multimodal Recordings during Educational Videos.
\newblock \emph{Scientific Data}, 13:920, 2026.

\end{thebibliography}


\end{document}
